% GENERATED FILE. Edit canonical lessons, then run npm run generate:books.
% canonical-source-hash: fnv1a64:23ceb4583d32538a
% canonical-lessons: TA-C541-kappitu, TA-C541-itar, TA-C541-apayam, TA-C541-mukavar, TA-C541-pinai

\chapter{Insurance and Risk}
\label{ch:ta-insurance-and-risk}
\hlchaptermodality{541}

\begin{chapteropening}
\textbf{By the end of this chapter:} \emph{I can talk about insurance and what happens if things go wrong.}

Complete the last lesson of chapter 541: I can talk about insurance and what happens if things go wrong.
\end{chapteropening}

\section[\texorpdfstring{\ta{காப்பீடு} (kāppīṭu)}{காப்பீடு (kāppīṭu)}]{\ta{காப்பீடு} (kāppīṭu) — insurance}

\label{lesson:TA-C541-kappitu}



\begin{quote}
Before the new one: say the Tamil for a rehearsal, then the Tamil for make-up.
\end{quote}

\subsection*{You'll want to know: \ta{காப்பீடு}}
\textbf{\ta{காப்பீடு}} — \emph{kāppīṭu} — \textquotedblleft{}insurance\textquotedblright{}.

\begin{quote}
\ta{வீட்டுக்குக்} \ta{காப்பீடு} \ta{இருந்தால்} \ta{கவலை} \ta{இல்லை}.
\end{quote}

\begin{quote}
\emph{vīṭṭukkuk kāppīṭu iruntāl kavalai illai.}
\end{quote}

\begin{quote}
If the house is insured, there is no worry.
\end{quote}

What is insured takes -\ta{க்கு}: \ta{வீட்டுக்குக்} \ta{காப்பீடு}. In speech people also say \ta{இன்சூரன்ஸ்}.

\begin{cousinweb}
\textbf{\ta{காப்பீடு}} (\emph{kāppīṭu}) is a modern native compound: \ta{காப்பு} (\emph{kāppu}), \textquotedblleft{}protection\textquotedblright{}, from \ta{கா} (\emph{kā}), \textquotedblleft{}to guard\textquotedblright{} (DEDR 1416), and a second part \emph{-īṭu}, whose exact sense here is not settled.
\end{cousinweb}

\begin{grammarlens}[title={What you've built}]
One more word for saying what happens if.
\end{grammarlens}

\subsection*{Your turn}
\begin{itemize}
\raggedright
  \item \emph{Say it:} \emph{kāppīṭu}
  \item \emph{Say it:} \emph{kāppīṭu}, once more
  \item \emph{Say it:} \emph{vīṭṭukkuk kāppīṭu iruntāl kavalai illai.} — If the house is insured, there is no worry.
  \item \emph{Recall:} say the Tamil for a rehearsal, then the Tamil for make-up, then say \emph{kāppīṭu} again
\end{itemize}

\begin{tcolorbox}[breakable,colback=teal!4,colframe=teal!35!black,title={Before you move on}]
What does \ta{காப்பீடு} mean? (\textquotedblleft{}Insurance\textquotedblright{}.) Say it once more.
\end{tcolorbox}
\section[\texorpdfstring{\ta{இடர்} (iṭar)}{இடர் (iṭar)}]{\ta{இடர்} (iṭar) — a risk; a hardship}

\label{lesson:TA-C541-itar}



\begin{quote}
Before the new one: say the Tamil for make-up, then the Tamil for insurance.
\end{quote}

\subsection*{You'll want to know: \ta{இடர்}}
\textbf{\ta{இடர்}} — \emph{iṭar} — \textquotedblleft{}a risk; a hardship\textquotedblright{}.

\begin{quote}
\ta{இந்த} \ta{முதலீட்டில்} \ta{இடர்} \ta{அதிகம்}.
\end{quote}

\begin{quote}
\emph{inta mutalīṭṭil iṭar atikam.}
\end{quote}

\begin{quote}
There is a lot of risk in this investment.
\end{quote}

A literary word; in insurance it means the risk covered.

\begin{cousinweb}
\textbf{\ta{இடர்}} (\emph{iṭar}) is native Dravidian: DEDR 435 gives \ta{இடர்} (\emph{iṭar}), \textquotedblleft{}affliction, distress, trouble\textquotedblright{}.
\end{cousinweb}

\begin{grammarlens}[title={What you've built}]
One more word for saying what happens if.
\end{grammarlens}

\subsection*{Your turn}
\begin{itemize}
\raggedright
  \item \emph{Say it:} \emph{iṭar}
  \item \emph{Say it:} \emph{iṭar}, once more
  \item \emph{Say it:} \emph{inta mutalīṭṭil iṭar atikam.} — There is a lot of risk in this investment.
  \item \emph{Recall:} say the Tamil for make-up, then the Tamil for insurance, then say \emph{iṭar} again
\end{itemize}

\begin{tcolorbox}[breakable,colback=teal!4,colframe=teal!35!black,title={Before you move on}]
What does \ta{இடர்} mean? (\textquotedblleft{}A risk; a hardship\textquotedblright{}.) Say it once more.
\end{tcolorbox}
\section[\texorpdfstring{\ta{அபாயம்} (apāyam)}{அபாயம் (apāyam)}]{\ta{அபாயம்} (apāyam) — danger; a hazard}

\label{lesson:TA-C541-apayam}



\begin{quote}
Before the new one: say the Tamil for insurance, then the Tamil for a risk.
\end{quote}

\subsection*{You'll want to know: \ta{அபாயம்}}
\textbf{\ta{அபாயம்}} — \emph{apāyam} — \textquotedblleft{}danger; a hazard\textquotedblright{}.

\begin{quote}
\ta{பாலம்} \ta{உடைந்தால்} \ta{அபாயம்} \ta{ஏற்படும்}.
\end{quote}

\begin{quote}
\emph{pālam uṭaintāl apāyam ēṟpaṭum.}
\end{quote}

\begin{quote}
If the bridge breaks, there will be danger.
\end{quote}

A danger sign says \ta{அபாயம்}. The adjective \ta{ஆபத்தான} (\emph{āpattāṉa}) belongs to a different word, \ta{ஆபத்து}.

\begin{cousinweb}
\textbf{\ta{அபாயம்}} (\emph{apāyam}) is borrowed from Sanskrit \emph{apāya}, \textquotedblleft{}departure, death, misfortune\textquotedblright{} (CDIAL 481), with the Tamil ending -\ta{அம்}.
\end{cousinweb}

\begin{grammarlens}[title={What you've built}]
One more word for saying what happens if.
\end{grammarlens}

\subsection*{Your turn}
\begin{itemize}
\raggedright
  \item \emph{Say it:} \emph{apāyam}
  \item \emph{Say it:} \emph{apāyam}, once more
  \item \emph{Say it:} \emph{pālam uṭaintāl apāyam ēṟpaṭum.} — If the bridge breaks, there will be danger.
  \item \emph{Recall:} say the Tamil for insurance, then the Tamil for a risk, then say \emph{apāyam} again
\end{itemize}

\begin{tcolorbox}[breakable,colback=teal!4,colframe=teal!35!black,title={Before you move on}]
What does \ta{அபாயம்} mean? (\textquotedblleft{}Danger; a hazard\textquotedblright{}.) Say it once more.
\end{tcolorbox}
\section[\texorpdfstring{\ta{முகவர்} (mukavar)}{முகவர் (mukavar)}]{\ta{முகவர்} (mukavar) — an agent}

\label{lesson:TA-C541-mukavar}



\begin{quote}
Before the new one: say the Tamil for a risk, then the Tamil for danger.
\end{quote}

\subsection*{You'll want to know: \ta{முகவர்}}
\textbf{\ta{முகவர்}} — \emph{mukavar} — \textquotedblleft{}an agent\textquotedblright{}.

\begin{quote}
\ta{காப்பீட்டு} \ta{முகவர்} \ta{நாளை} \ta{வீட்டுக்கு} \ta{வருவார்}.
\end{quote}

\begin{quote}
\emph{kāppīṭṭu mukavar nāḷai vīṭṭukku varuvār.}
\end{quote}

\begin{quote}
The insurance agent will come to the house tomorrow.
\end{quote}

One who acts for a company: \ta{பயண} \ta{முகவர்}, a travel agent.

\begin{cousinweb}
The origin of \textbf{\ta{முகவர்}} (\emph{mukavar}) is not settled: it looks built on \ta{முகம்} (\emph{mukam}), \textquotedblleft{}the face, the front\textquotedblright{} (DEDR 4889), with a personal ending, but no dictionary we can check shows the formation.
\end{cousinweb}

\begin{grammarlens}[title={What you've built}]
One more word for saying what happens if.
\end{grammarlens}

\subsection*{Your turn}
\begin{itemize}
\raggedright
  \item \emph{Say it:} \emph{mukavar}
  \item \emph{Say it:} \emph{mukavar}, once more
  \item \emph{Say it:} \emph{kāppīṭṭu mukavar nāḷai vīṭṭukku varuvār.} — The insurance agent will come to the house tomorrow.
  \item \emph{Recall:} say the Tamil for a risk, then the Tamil for danger, then say \emph{mukavar} again
\end{itemize}

\begin{tcolorbox}[breakable,colback=teal!4,colframe=teal!35!black,title={Before you move on}]
What does \ta{முகவர்} mean? (\textquotedblleft{}An agent\textquotedblright{}.) Say it once more.
\end{tcolorbox}
\section[\texorpdfstring{\ta{பிணை} (piṇai)}{பிணை (piṇai)}]{\ta{பிணை} (piṇai) — bail; a surety}

\label{lesson:TA-C541-pinai}



\begin{quote}
Before the new one: say the Tamil for danger, then the Tamil for an agent.
\end{quote}

\subsection*{You'll want to know: \ta{பிணை}}
\textbf{\ta{பிணை}} — \emph{piṇai} — \textquotedblleft{}bail; a surety\textquotedblright{}.

\begin{quote}
\ta{நீதிபதி} \ta{அவருக்குப்} \ta{பிணை} \ta{வழங்கினார்}.
\end{quote}

\begin{quote}
\emph{nītipati avarukkup piṇai vaḻaṅkiṉār.}
\end{quote}

\begin{quote}
The judge granted him bail.
\end{quote}

Also a person who stands surety for a loan. The verb \ta{பிணை} means \textquotedblleft{}to tie, to link\textquotedblright{}.

\begin{cousinweb}
\textbf{\ta{பிணை}} (\emph{piṇai}) is native Dravidian: DEDR 4160 gives \ta{பிணை} (\emph{piṇai}), \textquotedblleft{}to link, to tie, to fasten\textquotedblright{}, and the noun for being bound together; bail and surety are later senses.
\end{cousinweb}

\begin{grammarlens}[title={What you've built}]
One more word for saying what happens if.
\end{grammarlens}

\subsection*{Your turn}
\begin{itemize}
\raggedright
  \item \emph{Say it:} \emph{piṇai}
  \item \emph{Say it:} \emph{piṇai}, once more
  \item \emph{Say it:} \emph{nītipati avarukkup piṇai vaḻaṅkiṉār.} — The judge granted him bail.
  \item \emph{Recall:} say the Tamil for danger, then the Tamil for an agent, then say \emph{piṇai} again
\end{itemize}

\begin{tcolorbox}[breakable,colback=teal!4,colframe=teal!35!black,title={Before you move on}]
What does \ta{பிணை} mean? (\textquotedblleft{}Bail; a surety\textquotedblright{}.) Say it once more.
\end{tcolorbox}
